% Normative source: sections 2, 3, 5, 6, 14, 16.
\section{시스템 모델과 설계 목표}
\label{sec:system}

zkDPP는 원료의 등록·생산·이전을 회계 증표(Note)와 미해결 이전
증표(Voucher)의 생성·소비로 나타내고, 발급이 끝난 결과는 공개
제품 주장(DPP Claim)으로 남긴다. 이하 세 종류를 객체라고 한다.
\emph{소비}는 기록을 다음 거래의 입력으로 처리했다는 뜻이며,
물리적 원료가 사라졌다는 뜻은 아니다. 원장은 같은 기록의 재사용을
거부한다. 문제가 발견되면 감사자는 이 연결을 따라 아직 소비되지
않은 객체를 찾고 차단을 요청한다. 이 절에서는 정상 거래와 사후
대응을 누가 맡는지, 무엇을 신뢰하는지 정한다.

\subsection{참여 주체와 권한}

공급·생산·이전에 참여하는 기업은 원래 수치와 비밀 자료인
\emph{개방값}, 소유자임을 보일 \emph{소유권 비밀}을 보관하고
거래를 제출한다. 원장에 기록되는 각 거래를 이벤트(event)라고
부르며, Entry·Transfer 등은 그 종류를 나타낸다.
등록 승인자는 초기 값과 동일 물량의 중복 등록 여부를 판단한다.
정책 관리자는 가공·발급 조건과 검증키, 가공 권한을 등록한다.
원장의 검증 계약(router)은 제출된 증명과 현재 상태를 검사하고,
통과한 거래의 변경을 전부 적용하거나 전부 취소한다.
외부 DPP 검증자는 공개 문서와 확정된 발급 기록의 일치를 확인한다.

감사자와 상태 관리자는 같은 신뢰 주체로 둔다.
이 주체는 조사 시작점을 선택해 암호화된 연결을 읽고,
차단을 승인한 범위에서 남은 객체의 상태 변경을 요청한다.
감사 클라이언트는 그 주체의 소프트웨어이며 별도의 신뢰 주체가 아니다.

소유권은 공개 거래 주소와 별개다. 비밀 $sk$에 대해
$o=\mathsf{OwnerKey}(sk)$를 정의하고, 객체 소비의 증명에는
그 비밀이 필요하다. 소비 때 공개하는 \emph{소비 식별자}(nullifier)는
같은 객체를 다시 쓰지 못하게 한다. 소비 가능한 객체의 생성자는 이를 미리
계산해 감사 자료에 암호화한다. 이전 증표의 송신자도 자신이 만든
증표의 해결 식별자를 계산할 수 있다. 계산 방법과 이벤트별 소유권
규칙은 \ref{sec:accounting}절에서 설명한다.

\subsection{감사키와 접근 권한}

감사 자료는 공개키 $PK_B$로 암호화하고 비밀키 $SK_B$로 읽는다.
본문에서는 이를 $pk_{\mathsf{audit}}=PK_B$,
$sk_{\mathsf{audit}}=SK_B$로도 쓴다. 감사자가 승인 전에 자료를
읽지 못하도록 $SK_B$를 $N$개 비밀 조각으로 나누어 위원들이
하나씩 보관한다. 이는 단순히 키의 일부를 잘라 나누는 것이 아니라,
서로 다른 $K$개 조각으로 원래 키를 계산할 수 있게 하는
$K$-of-$N$ 비밀 분산이다($1\leq K\leq N$).
본 구성은 Shamir 방식으로 감사키 자체를 분산한다.~\cite{shamir-secret-sharing-1979}

초기 설정에서는 감사용 키쌍을 생성하고 비밀키를 분산한 뒤,
공개키와 구성원별 조각 확인 자료를 배포한다. 초기 설정자가
원본 키 사본을 남기지 않고 조각을 각 위원에게 비밀로 전달한다고
가정한다. 이 주체가 원본 키를 보유하면 위원들의 협력 없이도
감사 자료를 읽을 수 있으므로 초기 설정은 신뢰 경계에 포함된다.

외부 승인 절차를 통과하면 서로 다른 $K$명의 위원이 자신의
조각을 인증된 비밀 경로로 감사자에게 전달한다. 감사자는 각 조각의
구성원·배포·키 식별과 공개 확인 자료를 검사하고 $SK_B$를 복원한다.
복원한 키에서 계산한 공개키가 배포된 $PK_B$와 같을 때만
그 키를 사용한다. 따라서 별도 공개키로 감사키를 다시 암호화하는
단계는 필요하지 않다. 승인 판단과 조각 전달의 보안은 외부 절차의
책임이며, 이 논문은 그 절차 자체의 보안을 증명하지 않는다.

키를 복원한 감사자는 선택한 원장의 기록과 형식을 확인한 뒤,
기록별 암호문 $ct$를 자신의 컴퓨터에서 다음과 같이 복호화한다.
\begin{equation}
 M\leftarrow\mathsf{AuditDec}(sk_{\mathsf{audit}},ct)
 \label{eq:decrypt}
\end{equation}
키 식별·형식·복호화가 맞지 않으면 오류로 처리한다.
$SK_B$와 위원 조각은 공개 거래나 감사 결과에 넣지 않는다.
위원회는 감사키 복원에 협력하며, 기록마다 다시 참여할 필요는 없다.
이후 추적·차단 분석은 유효한 키를 복원한 감사자의 실행을 다룬다.

기본 실행은 배포마다 감사키 하나를 사용한다. 따라서 감사자는 한 번
$SK_B$를 복원하면 같은 키로 암호화된 다른 기록도 읽을 수 있다.
시작점·시점은 업무상 조사 범위일 뿐 복호화 범위를 제한하지 않는다.
키 교체나 승인 철회도 이미 전달한 키의 능력을 회수하지 못한다.

\subsection{신뢰 가정과 설계 목표}

참여자·증명자·DPP 게시자는 공모하거나 틀린 거래와 문서를 낼 수
있다. 반면 등록 승인자의 초기 판단, 정책 관리자의 등록,
감사자의 시작점 선택과 상태 집행은 신뢰한다. 비인가 공격자는
$SK_B$와 $K$개 이상의 위원 비밀 조각을 얻지 못한다고 가정한다.
초기 비밀 사본, 조각 전달 경로나 감사자에게서 키가 유출되는 경우는
이 가정 밖이다. 승인된 감사자에게는 감사 자료의 비밀성을
요구하지 않는다.

원장은 확정 순서·변경되지 않는 수락 기록·인증된 상태 읽기와
원자적 거래를 제공한다고 가정한다. 원격 조회 결과도 인증해야 한다.
이전 증표의 개방값은 송신자가 보관하고 인증된 비밀 채널로
수신자에게 전달한다. 전달 실패의 복구는 이 원장 설계의 보장이 아니다.

질량·재활용 크레딧·탄소와 운송·공정 증분은 외부에서 제공된다.
증명은 그 값의 범위와 배분 규칙을 검사하지만 측정의 진실성,
공정 수행이나 제품과 문서의 실제 일치는 확인하지 않는다.
탄소 배분 규칙 자체의 한계는 \ref{sec:limits}절에서 설명한다.

이 가정 아래 목표는 세 가지다. \emph{정확성}은 정상 조건의 요청이
기대한 결과를 내는 것, \emph{건전성}은 잘못된 거래·추적·완료
보고를 성공으로 받아들이지 않는 것이다. \emph{프라이버시}는
공개 정보가 같은 정상 이력에서 외부 관찰자가 숨긴 수치와 거래
연결을 구별하기 어렵게 하는 것이다. 차단 완료에는 소비 경합이
없는 등의 조건이 필요하고, 프라이버시는 참여자 비밀과 보호 대상의
감사·상태 정보가 공개되지 않은 경우로 한정한다. 정확한 정의와
증명은 \ref{sec:analysis}절에 제시한다.
